\documentclass{article}
\usepackage{enumitem}
\usepackage{amsmath}
\begin{document}

\title{Chapter 28: Molecular Basis of Inheritance}
\date{}
\maketitle

\begin{enumerate}[label=Q\arabic*.]

\item The Meselson--Stahl experiment proved semi-conservative replication using:
\\(A) Radioactive phosphorus ($^{32}$P) to label DNA
\\(B) Heavy nitrogen ($^{15}$N) to label DNA and CsCl density gradient centrifugation
\\(C) Tritium-labeled thymidine and autoradiography
\\(D) X-ray crystallography

\item DNA polymerase III in \textit{E. coli} synthesizes DNA:
\\(A) In both 3'$\rightarrow$5' and 5'$\rightarrow$3' directions
\\(B) Only in the 5'$\rightarrow$3' direction, adding nucleotides to the 3'-OH end
\\(C) In the 3'$\rightarrow$5' direction exclusively
\\(D) Without a template strand

\item The lac operon is an example of:
\\(A) Positive regulation only
\\(B) Negative regulation (repressor blocks transcription when lactose is absent) and positive regulation (CAP-cAMP when glucose is absent)
\\(C) Constitutive gene expression
\\(D) Attenuation

\item In eukaryotic gene expression, the 5' cap and poly-A tail on mRNA function to:
\\(A) Provide the start codon for translation
\\(B) Protect mRNA from exonuclease degradation and facilitate ribosome binding
\\(C) Act as transcription start signals
\\(D) Encode the signal peptide

\item Assertion: Satellite DNA is useful in DNA fingerprinting.\\
Reason: Satellite DNA consists of highly repetitive, non-coding sequences (STRs/VNTRs) that are highly variable between individuals.
\\(A) Both true, Reason is correct explanation
\\(B) Both true, Reason is not correct explanation
\\(C) Assertion true, Reason false
\\(D) Assertion false, Reason true

\item The human genome contains approximately how many protein-coding genes?
\\(A) 3 billion
\\(B) 100,000
\\(C) 20,000--25,000
\\(D) 3 million

\item RNA interference (RNAi) silences gene expression by:
\\(A) Blocking transcription of target gene
\\(B) Small interfering RNA (siRNA) guiding RISC complex to cleave complementary mRNA
\\(C) Methylation of target gene promoter
\\(D) Inhibiting translation at the ribosome directly

\item The template strand of DNA used for transcription is also called:
\\(A) Coding strand (sense strand)
\\(B) Non-template strand
\\(C) Antisense strand (template/non-coding strand)
\\(D) Leading strand

\item In translation, the start codon AUG codes for:
\\(A) Methionine (Met) in eukaryotes; N-formyl methionine (fMet) in prokaryotes
\\(B) Methionine in both prokaryotes and eukaryotes
\\(C) Valine in prokaryotes
\\(D) Any amino acid depending on context

\item The genetic code is described as ``degenerate'' because:
\\(A) Some codons do not code for any amino acid
\\(B) Most amino acids are encoded by more than one codon (synonymous codons)
\\(C) The same codon can code for different amino acids
\\(D) Codons can overlap

\item Splicing of pre-mRNA in eukaryotes involves removal of:
\\(A) Exons
\\(B) Introns by the spliceosome complex
\\(C) The 5' cap
\\(D) The poly-A tail

\item The Human Genome Project (HGP) revealed that:
\\(A) Humans have approximately 100,000 genes
\\(B) $\sim$98.5\% of the genome is non-coding; protein-coding genes are $\sim$1.5\%
\\(C) All DNA sequences are functional and code for proteins
\\(D) Genome size correlates perfectly with organismal complexity

\item Chromatin remodeling and histone acetylation generally:
\\(A) Condense chromatin, silencing gene expression
\\(B) Relax chromatin structure, allowing transcription factor access and gene activation
\\(C) Methylate DNA to silence genes
\\(D) Promote heterochromatin formation

\item Match the following enzymes in DNA replication with their functions:\\
1. DNA helicase $\rightarrow$ (a) Unwinds double helix by breaking H-bonds\\
2. Primase $\rightarrow$ (b) Synthesizes short RNA primer\\
3. DNA polymerase III $\rightarrow$ (c) Main enzyme for DNA synthesis (5'$\rightarrow$3')\\
4. DNA ligase $\rightarrow$ (d) Joins Okazaki fragments by sealing nicks
\\(A) 1-a, 2-b, 3-c, 4-d
\\(B) 1-b, 2-a, 3-d, 4-c
\\(C) 1-c, 2-d, 3-a, 4-b
\\(D) 1-d, 2-c, 3-b, 4-a

\item A mutation changes a codon from GAG (Glu) to GAA. This is an example of:
\\(A) Nonsense mutation
\\(B) Missense mutation
\\(C) Silent (synonymous) mutation (both code for Glu)
\\(D) Frameshift mutation

\end{enumerate}
\end{document}
